% Based on templates/lecture-example.tex.
% Original worked example: IO and Disk.pdf, slides 10.21--10.35.
\begingroup
\renewcommand{\recordid}[1]{}
\renewcommand{\relatedtopic}[2]{\hyperref[topic:#1]{#2}}
\exercisegroup{SC2005-L23-EX}{第 23 讲：磁盘调度路径与总移动量}

\textbf{日期与来源：}2026-11-06；IO and Disk 10.21--10.35（PDF pp.~21--35）。按讲义共同例题整理，补全路径与计数条件。

\linespread{1.15}\selectfont
\setlength{\parskip}{0.4em}

\exercisequestion{SC2005-L23-E01}{同一请求队列的五种调度结果}
\textbf{题意与来源。}讲义 10.21--10.35：磁盘柱面为 $0$--$199$，磁头从 $53$ 开始，等待队列按到达顺序为
\[
  98,\ 183,\ 37,\ 122,\ 14,\ 124,\ 65,\ 67.
\]
分别求 FCFS、SSTF、SCAN、C-SCAN、C-LOOK 的服务路径及总移动量。SCAN 初始向 $0$ 移动；C-SCAN、C-LOOK 向大编号服务。沿用讲义快照：不加入新请求，统计到本批最后一个请求完成，包含途中的边界和空程移动。对应\relatedtopic{SC2005-L23-T01}{磁盘调度}。

\begin{workedsolution}
\textbf{(1) FCFS，10.22--10.23。}直接按到达顺序：
\[
  53\to98\to183\to37\to122\to14\to124\to65\to67.
\]
\[
  D=45+85+146+85+108+110+59+2=\boxed{640}.
\]
顺序中多次跨越相距很远的柱面，移动量较大。

\textbf{(2) SSTF，10.24--10.26。}从 $53$ 看，$65$ 距离 $12$、$37$ 距离 $16$，先取 $65$；随后取距离 $2$ 的 $67$。从 $67$ 看，$37$ 距离 $30$，比 $98$ 的 $31$ 更近，故下一站为 $37$。逐次重算得到
\[
  53\to65\to67\to37\to14\to98\to122\to124\to183,
\]
\[
  D=12+2+30+23+84+24+2+59=\boxed{236}.
\]
每一步以\textbf{当前位置}重新比较剩余请求，不按它们距初始 $53$ 的距离一次排序。

\textbf{(3) SCAN，10.27--10.29。}先向小编号服务到边界 $0$，再反向：
\[
  53\to37\to14\to0\to65\to67\to98\to122\to124\to183.
\]
在每个单调移动段内，距离可以合并为端点差：
\[
  D=(53-0)+(183-0)=\boxed{236}.
\]
$0$ 虽无请求仍需经过；最后一个请求在 $183$，本题到此结束，不再增加到 $199$ 的路程。SCAN 与 SSTF 此次总量相等，但路径不同。

\textbf{(4) C-SCAN，10.30--10.32。}向大编号服务，到 $199$ 后空程回到 $0$，再向大编号服务：
\[
  53\to65\to67\to98\to122\to124\to183\to199\to0\to14\to37.
\]
\[
  D=(199-53)+(199-0)+(37-0)=146+199+37=\boxed{382}.
\]
$199\to0$ 不服务中间请求，但实际移动了 $199$ 个柱面间距。

\textbf{(5) C-LOOK，10.33--10.35。}只在待处理请求的范围内循环，最高为 $183$、最低为 $14$：
\[
  53\to65\to67\to98\to122\to124\to183\to14\to37.
\]
\[
  D=(183-53)+(183-14)+(37-14)=130+169+23=\boxed{322}.
\]
省去到 $199$ 和 $0$ 的移动，但 $183\to14$ 的空程仍需计入。

\textbf{LOOK 补充推导。}讲义 10.36 提到 LOOK，未对本队列展开计算。若保持 SCAN 的初始向小编号方向，但在最低请求 $14$ 处直接反向，则路径为
\[
  53\to37\to14\to65\to67\to98\to122\to124\to183,
\qquad D=(53-14)+(183-14)=208.
\]
这也给出 SSTF 并非全局最短路径的反例：同一请求集存在移动量 $208<236$ 的服务顺序。以上结果衡量的是柱面移动，换算时间还需设备模型。
\end{workedsolution}
\endgroup
